\documentclass[10pt,letterpaper]{article}
% ATS-friendly: one column, standard fonts, no tables, no images, no text boxes.
\usepackage[T1]{fontenc}
\usepackage[utf8]{inputenc}
\usepackage[scaled=0.92]{helvet}  % Arial-style vector font (Type 1), the safest choice for parsers
\renewcommand{\familydefault}{\sfdefault}
\usepackage{cmap}                 % makes the PDF text copy/paste and parse cleanly
\usepackage[letterpaper,left=0.65in,right=0.65in,top=0.45in,bottom=0.45in]{geometry}
\usepackage{microtype}
\DisableLigatures{encoding = *, family = *}   % "fi" and "fl" stay two plain letters for parsers
\usepackage{enumitem}
\usepackage[hidelinks]{hyperref}
\usepackage{titlesec}
\input{glyphtounicode}
\pdfgentounicode=1
\pagestyle{empty}
\setlength{\parindent}{0pt}
\setlength{\parskip}{0pt}
\raggedright
\titleformat{\section}{\large\bfseries}{}{0em}{}[\titlerule]
\titlespacing*{\section}{0pt}{6pt}{3pt}
\setlist[itemize]{leftmargin=1.2em,itemsep=0pt,topsep=1pt,parsep=0pt}
\begin{document}

\begin{center}{\LARGE\bfseries Piyush Chaudhary}\\[3pt]
\href{mailto:piyush.chaudhary@somaiya.edu}{piyush.chaudhary@somaiya.edu} \textbar{} \href{mailto:mail@piyushchy.is-a.dev}{mail@piyushchy.is-a.dev} \textbar{} \href{https://piyushchy.is-a.dev/}{piyushchy.is-a.dev} \textbar{} \href{https://github.com/Piyushchy}{github.com/Piyushchy} \textbar{} \href{https://www.linkedin.com/in/piyush-chaudhary-9b5999187/}{linkedin.com/in/piyush-chaudhary-9b5999187}
\end{center}

\section*{SUMMARY}
Final-year B.Tech Computer Engineering student who builds document-processing, cloud and mobile software. Currently an intern at SBI Global IT Centre, previously an intern at Mazagon Dock Shipbuilders Limited.

\section*{EDUCATION}
\noindent\textbf{Bachelor of Technology, Computer Engineering}, K.J. Somaiya College of Engineering\hfill Final year\par
Coursework: Natural Language Processing, Cloud Computing, AR/VR, Network and Information Security\par

\section*{EXPERIENCE}
\noindent\textbf{Intern}, State Bank of India, Global IT Centre\hfill Present\par
\begin{itemize}
  \item Building a production pipeline that extracts structured tables from scanned telecom audit PDFs into Excel, using a hybrid of computer vision (PyMuPDF, OpenCV) and a vision-language model.
  \item Compared three architectures on the same documents (Table Transformer, a pure vision-language model, and Docling) before choosing the hybrid design.
  \item Added a fast pre-scan that locates key header fields and the conclusion-table page range before the per-page pass, and switched to lazy page generators after eager rendering exhausted memory.
  \item Fixed a column-splicing bug that turned an 11-column table into 24 phantom columns, and added a guard that rejects implausible vision-language model cell text.
  \item Generated synthetic scanned PDFs with noise and skew for testing, and documented the system in a 16-page technical report.
\end{itemize}

\vspace{1pt}
\noindent\textbf{Intern}, Mazagon Dock Shipbuilders Limited\par

\section*{PROJECTS}
\noindent\textbf{The Tangled Path: Choice-Based Anti-Drug Game} (HTML, CSS, JavaScript)\hfill \href{https://sudo-choices.vercel.app/}{sudo-choices.vercel.app}\par
\begin{itemize}
  \item 1st runner up at a hackathon. Team leader and front-end developer on a team of four.
  \item A story game that follows one character through four life stages (ages 6, 15, 25 and 70), where each choice leads toward a good or a bad ending.
  \item Includes a built-in story teller and a full-screen mode. Deployed on Vercel.
\end{itemize}

\vspace{1pt}
\noindent\textbf{MouseShare} (Python, Kotlin, Android)\par
\begin{itemize}
  \item Built because I could not find an existing tool that shares a Windows PC's mouse and keyboard with an Android phone: a Python desktop app captures input and a Kotlin Android app receives and acts on it.
  \item Comparing transports (HID keyboard and mouse, WiFi) to cut input lag.
\end{itemize}

\vspace{1pt}
\noindent\textbf{GIF Crop and Overlay Tool} (JavaScript, HTML, Canvas)\hfill \href{https://piyushchy.github.io/gif-tool/}{piyushchy.github.io/gif-tool}\par
\begin{itemize}
  \item Built because I could not find an existing tool that crops a GIF to a freehand or polygon shape or overlays one GIF on another with frame-level control.
  \item Overlay mode has two frame tracks to slide against each other, plus per-frame blank, skip and hold. Wrote the GIF decoder and encoder from scratch (LZW, median-cut palette). Runs entirely in the browser.
\end{itemize}

\vspace{1pt}
\noindent\textbf{Cloud-Native App on AWS} (Docker, AWS ECR, ECS Fargate, ALB, RDS MySQL, CodePipeline, SNS)\par
\begin{itemize}
  \item Containerized a Node.js app and deployed it on ECS Fargate behind an Application Load Balancer, with the image stored in ECR.
  \item Connected to an RDS MySQL database through a bastion host over Session Manager, and added a CodePipeline CI/CD flow and an SNS alarm.
\end{itemize}

\vspace{1pt}
\noindent\textbf{RNN vs LSTM Text Classification} (Python, PyTorch, Google Colab)\par
\begin{itemize}
  \item Trained and compared a plain RNN and an LSTM on the same text classification task in a runnable Google Colab notebook.
\end{itemize}

\section*{TECHNICAL SKILLS}
\textbf{Languages:} Python, Java, Kotlin, JavaScript, HTML, CSS\par
\textbf{Web:} React.js, HTML5, CSS3\par
\textbf{Computer vision and document AI:} OpenCV, PyMuPDF, EasyOCR, Table Transformer, Docling, Pillow, vision-language models\par
\textbf{Machine learning:} PyTorch, RNN, LSTM\par
\textbf{Cloud and infrastructure:} Docker, AWS (ECR, ECS Fargate, ALB, RDS MySQL, EC2, CodePipeline, SNS)\par
\textbf{Tools and systems:} Linux, Command Line (CLI), SSH, Git, GitHub, Google Colab, openpyxl\par
\end{document}
